\section{Version Control Tools}

El control de versiones es una práctica esencial en el desarrollo de software basada en gestionar y almacenar todos los cambios realizados a lo largo del tiempo sobre un producto.
Aunque un control de versiones puede realizarse de forma manual, es prácticamente imperativo el hecho de utilizar un sistema que lo administre automáticamente.

\vspace{10px}

Dos de los más conocidos son Apache Subversion (SVN) y Git, por lo que nos centraremos en estos dos para discernir cual de ellos es preferible integrar en nuestra \textit{pipeline} de desarrollo.

\subsection{SVN}
Subversion (SVN) comenzó su desarrollo a principios del año 2000 por \textit{CollabNet}, hasta que, hacia el 2009, el proyecto fue absorbido por la \textit{Apache Software Foundation}.

\vspace{10px}

SVN se basa en un sistema de control de versiones centralizado, por lo que existirá un único repositorio visible y accesible para todos los usuarios. Este hecho hace imposible fusionar cambios, por lo que un archivo no puede ser editado por más de un usuario de forma concurrente.

\subsection{Git}
En 2005, Linus Torvalds, el creador de Linux, probablemente el kernel monolítico de código abierto basado en Unix más usado en el mundo, inició el desarrollo de un nuevo sistema de control de versiones como alternativa a Bitkeeper.

\vspace{10px}

Git se basa en un modelo distribuido, por lo que todos los usuarios disponen del repositorio de forma local. Para permitir que todos los cambios se unifiquen, todos los repositorios locales se deberán sincronizar con uno o varios repositorios remotos, donde se incorporan y unen todos los cambios.
De esta forma, todos los desarrolladores tienen acceso al repositorio completo, incluyendo el historial local, sin depender de ningún tipo de conexión de red. Cuando alguien tenga listo un cambio en local, se puede cargar con una gran velocidad al repositorio remoto, sin ningún tipo de bloqueo.

\pagebreak
